\section{Per-search options are not additional GUCs}
\label{sec:appendix-search-options}
\label{sec:appendix-options-precedence}

The third argument of \texttt{mmsearch} accepts the
\nolinkurl{bioseq_aligner_options} composite. The
\texttt{aligner\_options(...)} constructor creates that value. Both names
belong to the extension's installation schema, assumed to be
\texttt{public} in the example below. The sequence type remains in
\texttt{bioseq}.

All 36 constructor arguments default to SQL \texttt{NULL}. They do not
capture the current GUC values when the composite is constructed.
When the search resolves its parameters, each non-NULL field replaces its
corresponding inherited value. An omitted field and an explicitly NULL
field both leave that default unchanged. Zero and \texttt{false} are
explicit overrides. The independent \texttt{index\_read\_requirement}
starts at \texttt{allow}, encoded as 0.

\begin{lstlisting}[style=bioseq-search-sql,
  caption={Construct per-search overrides without changing the session settings.},
  label={lst:configuration-options}]
SELECT (public.aligner_options(
  minimizer_algorithm => 0,
  index_read_requirement => 1)).*;
\end{lstlisting}

Here code 0 selects scalar minimizers and code 1 disables persisted
target-index reads for that search. The constructor alone performs no
sequence search. The same composite can be passed to an eligible stored
column query or a sequence join.

The options must resolve during planning to a constant of the exact
composite type. A constructor with constant arguments satisfies this rule.
Row-dependent options and unresolved generic-plan parameters are not
supported. Omitting the third argument, using an empty constructor and
passing \nolinkurl{NULL::public.bioseq_aligner_options} all request inherited
options.

Table~\ref{tab:configuration-options} gives every field in constructor order.
The first 35 fields map to GUCs documented above. The table's
\texttt{double} type means SQL \texttt{double precision}. These four fields
are converted to single precision in the alignment parameters.
The constructor only forms the composite value. The planner's decoder
validates the algorithm and index-requirement codes when a search uses it.
It does not enforce the ordinary numeric GUC intervals on the other
per-call fields. Kernel requirements still apply.

% Generated by generate_reference.py.
\begingroup
\footnotesize
\setlength{\tabcolsep}{3pt}
\renewcommand{\arraystretch}{1.12}
\begin{longtable}{@{}>{\raggedright\arraybackslash}p{.28\linewidth}>{\raggedright\arraybackslash}p{.10\linewidth}>{\raggedright\arraybackslash}p{\dimexpr.62\linewidth-4\tabcolsep\relax}@{}}
\caption{All nullable alignment options (for \texttt{aligner\_options} call) in the order as specified in the SQL API}\label{tab:configuration-options}\\
\toprule
\textbf{Options field} & \textbf{SQL type} & \textbf{Inherited GUC or independent policy}\\
\midrule
\endfirsthead
\toprule
\textbf{Options field} & \textbf{SQL type} & \textbf{Inherited GUC or independent policy}\\
\midrule
\endhead
\midrule
\multicolumn{3}{r}{\textit{Continued on the next page}}\\
\endfoot
\bottomrule
\endlastfoot
\nolinkurl{kmer_size} & int & \nolinkurl{bioseq.aligner.minimizer.kmer_size}\\[3pt]
\nolinkurl{window_size} & int & \nolinkurl{bioseq.aligner.minimizer.window_size}\\[3pt]
\nolinkurl{max_gap} & int & \nolinkurl{bioseq.aligner.chaining.maximum_gap}\\[3pt]
\nolinkurl{max_gap_ref} & int & \nolinkurl{bioseq.aligner.chaining.maximum_reference_gap}\\[3pt]
\nolinkurl{bandwidth} & int & \nolinkurl{bioseq.aligner.chaining.bandwidth}\\[3pt]
\nolinkurl{bandwidth_long} & int & \nolinkurl{bioseq.aligner.chaining.long_join_bandwidth}\\[3pt]
\nolinkurl{max_chain_skip} & int & \nolinkurl{bioseq.aligner.chaining.maximum_skipped_anchors}\\[3pt]
\nolinkurl{max_chain_iter} & int & \nolinkurl{bioseq.aligner.chaining.maximum_predecessors}\\[3pt]
\nolinkurl{min_seed_count} & int & \nolinkurl{bioseq.aligner.chaining.minimum_seed_count}\\[3pt]
\nolinkurl{min_chain_score} & int & \nolinkurl{bioseq.aligner.chaining.minimum_score}\\[3pt]
\nolinkurl{chain_gap_scale} & double & \nolinkurl{bioseq.aligner.chaining.gap_penalty_scale}\\[3pt]
\nolinkurl{chain_skip_scale} & double & \nolinkurl{bioseq.aligner.chaining.skip_penalty_scale}\\[3pt]
\nolinkurl{mask_level} & double & \nolinkurl{bioseq.aligner.secondary.overlap_threshold}\\[3pt]
\nolinkurl{mask_len} & int & \nolinkurl{bioseq.aligner.secondary.uncovered_length_limit}\\[3pt]
\nolinkurl{pri_ratio} & double & \nolinkurl{bioseq.aligner.secondary.score_ratio_threshold}\\[3pt]
\nolinkurl{best_n} & int & \nolinkurl{bioseq.aligner.secondary.maximum_count}\\[3pt]
\nolinkurl{match_score} & int & \nolinkurl{bioseq.aligner.scoring.match_score}\\[3pt]
\nolinkurl{mismatch_pen} & int & \nolinkurl{bioseq.aligner.scoring.mismatch_penalty}\\[3pt]
\nolinkurl{gap_open1} & int & \nolinkurl{bioseq.aligner.scoring.short_gap_open_penalty}\\[3pt]
\nolinkurl{gap_ext1} & int & \nolinkurl{bioseq.aligner.scoring.short_gap_extension_penalty}\\[3pt]
\nolinkurl{gap_open2} & int & \nolinkurl{bioseq.aligner.scoring.long_gap_open_penalty}\\[3pt]
\nolinkurl{gap_ext2} & int & \nolinkurl{bioseq.aligner.scoring.long_gap_extension_penalty}\\[3pt]
\nolinkurl{sc_ambi} & int & \nolinkurl{bioseq.aligner.scoring.ambiguous_base_penalty}\\[3pt]
\nolinkurl{zdrop} & int & \nolinkurl{bioseq.aligner.scoring.z_drop_threshold}\\[3pt]
\nolinkurl{end_bonus} & int & \nolinkurl{bioseq.aligner.scoring.end_bonus}\\[3pt]
\nolinkurl{max_occ} & int & \nolinkurl{bioseq.aligner.seeding.maximum_occurrence}\\[3pt]
\nolinkurl{mid_occ} & int & \nolinkurl{bioseq.aligner.seeding.intermediate_occurrence_threshold}\\[3pt]
\nolinkurl{block_min_size} & int & \nolinkurl{bioseq.aligner.block_alignment.minimum_block_size}\\[3pt]
\nolinkurl{block_max_size} & int & \nolinkurl{bioseq.aligner.block_alignment.maximum_block_size}\\[3pt]
\nolinkurl{chaining_algorithm} & int & \nolinkurl{bioseq.aligner.chaining.algorithm}\par Integer codes should be checked in the GUC table.\\[3pt]
\nolinkurl{base_alignment_algorithm} & int & \nolinkurl{bioseq.aligner.base_alignment.algorithm}\par Integer codes should be checked in the GUC table.\\[3pt]
\nolinkurl{minimizer_algorithm} & int & \nolinkurl{bioseq.aligner.minimizer.algorithm}\par Integer codes should be checked in the GUC table.\\[3pt]
\nolinkurl{vulkan_enabled} & bool & \nolinkurl{bioseq.aligner.vulkan.enabled}\\[3pt]
\nolinkurl{vulkan_seed_hits_enabled} & bool & \nolinkurl{bioseq.aligner.vulkan.seed_lookup.enabled}\\[3pt]
\nolinkurl{vulkan_index_cache_enabled} & bool & \nolinkurl{bioseq.aligner.vulkan.index_cache.enabled}\\[3pt]
\nolinkurl{index_read_requirement} & int & No GUC. See Table~\ref{tab:configuration-index-requirement}.\\[3pt]
\end{longtable}
\endgroup


\subsection{The index-read requirement}
\label{sec:appendix-options-index-policy}

This option has no GUC equivalent. It describes what a particular search
requires, independently of how its target index was built.

\begingroup
\small
\setlength{\tabcolsep}{3pt}
\begin{longtable}{@{}>{\raggedright\arraybackslash}p{.25\linewidth}
  >{\raggedright\arraybackslash}p{.07\linewidth}
  >{\raggedright\arraybackslash}p{.25\linewidth}
  >{\raggedright\arraybackslash}p{\dimexpr.43\linewidth-6\tabcolsep\relax}@{}}
\caption{The per-search index policy, separate from the GUC inventory.}
\label{tab:configuration-index-requirement}\\
\toprule
\textbf{Options field} & \textbf{Type} & \textbf{Value} & \textbf{Description}\\
\midrule
\endfirsthead
\toprule
\textbf{Options field} & \textbf{Type} & \textbf{Value} & \textbf{Description}\\
\midrule
\endhead
\bottomrule
\endfoot
\multirow[t]{3}{=}{\nolinkurl{index_read_requirement}} &
\multirow[t]{3}{=}{int} &
\textbf{0: \texttt{allow}}\par(default) &
Controls whether the search may depend on a persisted target index.
Follow the optional-read GUC and use a compatible supported reader when
available. The executor may instead compute a temporary lookup.\\*
& & 1: \nolinkurl{require_absent} &
Execute without persisted-index reads or a reused assembled persisted
lookup. Existing indexes are not deleted. Separate eligible background
persistence is not prohibited.\\*
& & 2: \nolinkurl{require_compatible_persisted} &
Require a compatible saved generation, overriding the optional-read GUC.
Missing or incompatible storage is an error. Recomputing the target index
cannot satisfy this requirement.\\
\end{longtable}
\endgroup

The prepared persisted reader belongs to an eligible two-input
\texttt{BioseqBatchJoin}. An ordinary executable \texttt{BioseqScan}
rejects requirement 2. A transformed target, such as a slice or
concatenation, must use its evaluated sequence view. That path forces
requirement 1 and rejects requirement 2.

\subsection{Overrides and execution requirements}
\label{sec:appendix-options-algorithms}

The main GUC tables show integer codes for all three algorithm fields.
Their inherited built-in choices are minimizer code 4
(\texttt{bounded\_nthash\_simd}), chaining code 1
(\texttt{simd\_minimap2}) and base-alignment code 1
(\texttt{block\_aligner}). These are inherited choices, not non-NULL
constructor defaults.

Persisted readers must still prepare query minimizers compatible with their
saved target profile. An accepted per-search override does not guarantee
that it can use an existing index.

An explicit \texttt{vulkan\_enabled => true} also requests a standalone
Vulkan prefilter in an eligible ordinary scan with an exact edit-distance
filter. Merely inheriting the enabled GUC does not make that extra request.
Consequently, an explicit per-call Boolean and the corresponding session
default are not interchangeable in that case.

Base alignment runs when the requested result requires it. Choosing
\nolinkurl{base_alignment_algorithm} alone does not cause a coordinate-only
search to calculate a base-alignment score. Requests for CIGAR,
\texttt{has\_cigar} or edit-related results select the additional work.
BlockAligner's first gap penalties
must be positive with a sum no greater than 128. Effective block sizes must
be powers of two with the minimum no greater than the maximum. The GUC
numeric ranges do not express all these joint conditions.
